% Tabasco descriptor-regression probe (R^2) -- \input{} from report-draft.tex Ch5.
\begin{table}[tbp]
\centering
\small
\begin{tabular}{l c c c c}
\toprule
\textbf{Source} & \textbf{MolWt} & \textbf{LogP} & \textbf{\#Rings} & \textbf{\#RotBonds} \\
\midrule
\multicolumn{5}{l}{\textit{Frozen encoders (the REPA targets in isolation)}} \\
CheMeleon (2D) & 0.993 & 0.988 & 0.995 & 0.989 \\
MACE (3D)      & 0.022 & 0.307 & 0.081 & $-0.014$ \\
Dummy          & 0.822 & 0.160 & 0.432 & 0.345 \\
\midrule
\multicolumn{5}{l}{\textit{Generator trunk (the model's internal representation)}} \\
Baseline             & 0.920 & 0.713 & 0.742 & 0.636 \\
REPA-CheMeleon       & 0.943 & 0.669 & 0.736 & 0.606 \\
REPA-MACE            & 0.935 & 0.703 & 0.704 & 0.582 \\
REPA-MACE (tradeoff) & 0.955 & 0.730 & 0.732 & 0.666 \\
\bottomrule
\end{tabular}
\caption{\textbf{Linear-probe $R^2$ reveals little representational difference between REPA and the baseline.} The baseline trunk already captures what CheMeleon learns, while MACE's embedding is too narrow to decode these descriptors. Hence, neither makes for a usable REPA target. ($R^2$ for four RDKit descriptors, on held-out molecules; \S\ref{sec:tabasco-metrics}; $n{=}1$ seed.)}
\label{tab:tabasco-probe}
\end{table}
